% Rascunho da seção de metodologia (abordagem emulada). Texto final com a equipe.
% Números vêm de % Gerado por src/analysis/artigo_tex.py a partir de results/. Não edite à mão.
\newcommand{\bandaQuatroCinco}{1,7}
\newcommand{\bandaQuatroQuinze}{5,2}
\newcommand{\bandaQuatroTrinta}{10,5}
\newcommand{\bandaSeisCinco}{3,0}
\newcommand{\bandaSeisQuinze}{8,9}
\newcommand{\bandaSeisTrinta}{17,7}
\newcommand{\bandaTresCinco}{0,8}
\newcommand{\bandaTresQuinze}{2,5}
\newcommand{\bandaTresTrinta}{5,0}
\newcommand{\etoeSsdFpQuatroNuvem}{102}
\newcommand{\etoeSsdFpTresNuvem}{87}
\newcommand{\etoeYoloFpQuatroNuvem}{111}
\newcommand{\etoeYoloFpSeisNuvem}{165}
\newcommand{\etoeYoloFpTresNuvem}{77}
\newcommand{\mapSsdFpQuatroNuvem}{0,595}
\newcommand{\mapSsdFpTresNuvem}{0,592}
\newcommand{\mapYoloFpQuatroNuvem}{0,745}
\newcommand{\mapYoloFpSeisNuvem}{0,761}
\newcommand{\mapYoloFpTresNuvem}{0,668}
\newcommand{\nCenarios}{15}
\newcommand{\nCombBorda}{0}
\newcommand{\nCombNuvem}{5}
\newcommand{\nQuadrosValidos}{300}
 (gerado por src/analysis/artigo_tex.py); não digite valores à mão.
\section{Metodologia}
\label{sec:metodologia}

\subsection{Caracterização}
Trata-se de uma pesquisa experimental conduzida por emulação. Sem acesso ao hardware físico
(Raspberry Pi 4/5 e acelerador Coral Edge TPU), o dispositivo de borda é representado por um
contêiner Docker \texttt{linux/arm64} executado no x86-64 por meio do emulador QEMU, limitado a
4 núcleos e 4\,GB de memória (perfil do Raspberry Pi~4B) e com o modelo de processador
Cortex-A72 (\texttt{QEMU\_CPU=cortex-a72}). Todo o código, as versões fixadas das ferramentas e os
dados brutos estão no repositório do experimento, e cada etapa é refeita com um único comando.

\subsection{Modelos e quantização}
Foram avaliados três detectores leves pré-treinados no COCO: YOLOv8n (Ultralytics), SSD-MobileNet~V1
(TF1 Detection Model Zoo) e EfficientDet-Lite0 (MediaPipe), cada um em FP32 e INT8, no formato
TensorFlow Lite (LiteRT). O INT8 do SSD-MobileNet~V1 e do EfficientDet-Lite0 é inteiro total, obtido por
quantização pós-treinamento calibrada nas mesmas 128 imagens (COCO128); o INT8 do YOLOv8n é o
export oficial da Ultralytics, que mantém entrada, saída e cabeça de detecção em float32 (INT8
híbrido). Os modelos INT8 também foram submetidos ao Edge TPU Compiler (versão 14.1), cujo relatório
indica quais operações seriam executadas no acelerador (Tabela~\ref{tab:edgetpu}).

\subsection{Fatores e amostragem}
Os fatores são o modelo (3 níveis), a precisão (FP32 e INT8), a resolução do quadro da fonte (320, 480
e 640 pixels) e a taxa da fonte (5, 15 e 30 FPS), totalizando 54 cenários na borda. A resolução é a do
quadro entregue pela câmera: o YOLOv8n usa o modelo exportado na mesma resolução, enquanto o
SSD-MobileNet~V1 (entrada fixa de $300\times300$) e o EfficientDet-Lite0 ($320\times320$) redimensionam o
quadro, e esse custo entra no pré-processamento.

Por restrição de tempo de emulação, o desenho original (1.050 inferências em cada cenário) foi
reduzido por decisão da equipe: cada uma das 18 combinações modelo $\times$ precisão $\times$ resolução
foi medida uma única vez, em processo novo, com 350 quadros, dos quais os 50 primeiros foram
descartados como aquecimento (\nQuadrosValidos{} quadros válidos). Como a taxa da fonte não altera o
tempo de processamento de cada quadro, os cenários de 5, 15 e 30 FPS foram derivados dessas
medições: o quadro $k$ fica disponível em $k/\text{fps}$, o processamento começa quando o quadro está
disponível e o anterior terminou, e a espera é registrada como tempo de fila, separada da latência
fim a fim. Também se contabilizam os quadros que uma política de descarte (processar sempre o
quadro mais recente) deixaria de processar.

\subsection{Conjunto de teste e qualidade}
O conjunto de teste tem 350 imagens do COCO val2017 com pessoas, sem anotações de multidão,
balanceadas entre 1, 2 e 3 ou mais pessoas e ordenadas pelo identificador; a lista e o sha256 de cada
imagem estão no repositório. Cada imagem é redimensionada para a resolução da fonte e entregue como
quadro. A qualidade é medida pelo mAP@50 da classe pessoa (pycocotools) e pelo F1 com limiar de
pontuação 0,5 e IoU 0,5, a partir das detecções salvas na própria coleta. A decodificação das saídas
e a supressão de não máximos (IoU 0,5, até 100 detecções) são feitas em Python com os mesmos
parâmetros para os três modelos, e foram validadas no x86 antes da coleta.

\subsection{Métricas de desempenho}
Para cada quadro são registrados o tempo de pré-processamento, de inferência
(\texttt{Interpreter.invoke}, 4 threads) e de pós-processamento, cuja soma é a latência fim a fim, além
do uso de CPU do processo e da memória residente. A taxa sustentável é o inverso da latência média.
A energia por quadro seria estimada pelo produto entre a potência de referência da placa e a latência,
mas não foi calculada por falta de um valor de potência com fonte verificável.

\subsection{Baseline de nuvem}
Para a QP2, o mesmo conjunto de quadros é enviado a um servidor de inferência: dois contêineres
\texttt{linux/amd64} em uma rede Docker, com o enlace limitado por \texttt{tc netem} a um RTT de 40\,ms e
subida de 10\,Mbit/s (valores assumidos). O cliente apenas codifica cada quadro em JPEG (qualidade 80)
e o envia por HTTP; o servidor executa os mesmos modelos FP32. São registrados os tempos de
codificação, transmissão, inferência e retorno, e o tamanho de cada quadro, de onde se calcula a banda
de subida necessária em cada taxa da fonte. Por simplicidade e para manter o pareamento por quadro,
usou-se envio quadro a quadro em JPEG no lugar de RTSP ou WebRTC.

\subsection{Análise estatística}
Para cada cenário calculam-se média, desvio padrão, mediana, percentil 95 e intervalo de confiança de
95\% da média. As comparações FP32 $\times$ INT8 e borda $\times$ nuvem são pareadas por quadro: aplica-se
o teste de Shapiro-Wilk às diferenças e, conforme o resultado, o teste $t$ pareado ou o de Wilcoxon,
com $\alpha = 0{,}05$.

\subsection{Limitações}
Os tempos de borda vêm de emulação e não representam a velocidade de um Raspberry Pi: servem para
comparações relativas (FP32 $\times$ INT8, entre modelos e resoluções). Além disso, sob o QEMU a
biblioteca Python do LiteRT mostrou-se cerca de dez vezes mais lenta que a ferramenta nativa
\texttt{benchmark\_model}. A comparação absoluta entre a latência da borda emulada e a da nuvem nativa
não é válida; o resultado da QP2 é a banda de subida. Os cenários de taxa da fonte são derivados, a
rede da nuvem usa valores assumidos, não há gravações residenciais e o Edge TPU foi avaliado apenas
pelo relatório do compilador.
